%
% 6.006 problem set 1 solutions template
%
\documentclass[12pt,twoside]{article}

\input{macros-sp20}
\newcommand{\theproblemsetnum}{1}

\title{6.006 Problem Set 1}

\begin{document}

\handout{Problem Set \theproblemsetnum}

\setlength{\parindent}{0pt}
\medskip\hrulefill\medskip

{\bf Name:} Your Name

\medskip

{\bf Collaborators:} Name1, Name2

\medskip\hrulefill

%%%%%%%%%%%%%%%%%%%%%%%%%%%%%%%%%%%%%%%%%%%%%%%%%%%%%
% See below for common and useful latex constructs. %
%%%%%%%%%%%%%%%%%%%%%%%%%%%%%%%%%%%%%%%%%%%%%%%%%%%%%

% Some useful commands:
%$f(x) = \Theta(x)$
%$T(x, y) \leq \log(x) + 2^y + \binom{2n}{n}$
% {\tt code\_function}


% You can create unnumbered lists as follows:
%\begin{itemize}
%    \item First item in a list
%        \begin{itemize}
%            \item First item in a list
%                \begin{itemize}
%                    \item First item in a list
%                    \item Second item in a list
%                \end{itemize}
%            \item Second item in a list
%        \end{itemize}
%    \item Second item in a list
%\end{itemize}

% You can create numbered lists as follows:
%\begin{enumerate}
%    \item First item in a list
%    \item Second item in a list
%    \item Third item in a list
%\end{enumerate}

% You can write aligned equations as follows:
%\begin{align}
%    \begin{split}
%        (x+y)^3 &= (x+y)^2(x+y) \\
%                &= (x^2+2xy+y^2)(x+y) \\
%                &= (x^3+2x^2y+xy^2) + (x^2y+2xy^2+y^3) \\
%                &= x^3+3x^2y+3xy^2+y^3
%    \end{split}
%\end{align}

% You can create grids/matrices as follows:
%\begin{align}
%    A =
%    \begin{bmatrix}
%        A_{11} & A_{21} \\
%        A_{21} & A_{22}
%    \end{bmatrix}
%\end{align}

% You can include images and PDFs as follows:
% \includegraphics[width=0.5\textwidth]{img.jpg}

\begin{problems}

\problem  % Problem 1

\begin{problemparts}

\problempart % Problem 1a
$(f_5, f_3, f_4, f_1, f_2)$

Compare growth by looking at $\ln$ of each function. $f_5 = \log\log(6006n) = \Theta(\log\log n)$
barely grows, so it's smallest. $f_3 = \log(n^{6006}) = 6006\log n = \Theta(\log n)$, and $\log n$
dominates $\log\log n$, so $f_5$ comes before $f_3$. $f_4 = (\log n)^{6006}$ is a fixed power of
$\log n$, which eventually beats any constant multiple of $\log n$, so $f_3$ comes before $f_4$.
For $f_4$ vs.\ $f_1 = \log(n^n) = n\log n$: $\ln f_4 = 6006\ln\ln n$ while $\ln f_1 \approx \ln n$,
and $\ln n$ eventually dominates $\ln\ln n$, so $f_4$ is $o(f_1)$. Finally $f_2 = (\log n)^n$ beats
$f_1$ since $\ln f_2 = n\ln\ln n$ grows faster than $\ln f_1 \approx \ln n$.

\problempart % Problem 1b
$(f_1, f_2, f_5, f_4, f_3)$

Write everything as a power of 2: $f_1 = 2^n$, $f_2 = 6006^n = 2^{n\log_2 6006}$,
$f_5 = 6006^{n^2} = 2^{n^2\log_2 6006}$, $f_4 = 6006^{2^n} = 2^{2^n\log_2 6006}$, and
$f_3 = 2^{6006^n} = 2^{2^{n\log_2 6006}}$. Since $\log_2 6006 > 1$, the exponents satisfy
$n < n\log_2 6006 < n^2\log_2 6006 < 2^n\log_2 6006 < 2^{n\log_2 6006}$ for large $n$
(each step is either multiplying by a growing factor or swapping a polynomial exponent for
an exponential one), which gives exactly this order.

\problempart % Problem 1c
$(\{f_2, f_5\}, f_4, f_1, f_3)$

$f_2 = \binom{n}{n-6} = \binom{n}{6} = \Theta(n^6)$, which ties with $f_5 = n^6$. By the
entropy bound on binomial coefficients, $f_4 = \binom{n}{n/6} = 2^{\Theta(n)}$ --- exponential
with a linear exponent --- which dominates any polynomial, so $\{f_2,f_5\}$ comes before $f_4$.
Next, $f_1 = n^n = 2^{n\log n}$ grows faster than any $2^{cn}$, so $f_4$ comes before $f_1$.
Finally, Stirling's approximation gives $\ln f_3 = \ln((6n)!) \approx 6n\ln(6n) - 6n \approx
6n\ln n$, versus $\ln f_1 \approx n\ln n$; since the gap $5n\ln n \to \infty$, $f_3$ dominates $f_1$.

\problempart % Problem 1d
$(f_5, f_2, f_1, f_3, f_4)$

$f_5 = n^{12+1/n} = \Theta(n^{12})$ is the only true polynomial here, so it's smallest.
Writing the rest as powers of 2: by Stirling, $n! = o(n^{n+4})$, so
$f_1 = n^{n+4}+n! = \Theta(n^{n+4}) = 2^{\Theta(n\log n)}$. Also $f_2 = n^{7\sqrt n} =
2^{7\sqrt n \log n}$, $f_3 = 4^{3n\log n} = 2^{6n\log n}$, and $f_4 = 7^{n^2} =
2^{n^2\log_2 7}$. Since $\sqrt n \log n = o(n\log n)$, $f_2$ precedes $f_1$. Since the exponent
gap between $f_3$ and $f_1$ is $\Theta(n\log n)\to\infty$, $f_1$ precedes $f_3$. And since
$n\log n = o(n^2)$, $f_3$ precedes $f_4$.

\end{problemparts}

\newpage
\problem  % Problem 2

\begin{problemparts}
\problempart % Problem 2a
To reverse the $k$ items starting at index $i$, pull them out one at a time from the front of
the block (which shifts the rest of the block down by one each time), then reinsert them at
the same starting index in the opposite order:

\begin{verbatim}
def reverse(D, i, k):
    items = []
    for t in range(k):
        items.append(D.delete_at(i))      # extracts items in order
    for t in range(k):
        D.insert_at(i + t, items[k - 1 - t])  # writes them back reversed
\end{verbatim}

Each of the $2k$ calls to \texttt{delete\_at}/\texttt{insert\_at} costs $O(\log n)$, so the
total runtime is $O(k\log n)$.

\problempart % Problem 2b
Pull the $k$ items out of position $i$ (again, $k$ deletions at the same index $i$, since the
block shifts down after each one), and reinsert them, in their original order, in front of
item $j$. The only subtlety is that if $j \geq i+k$, deleting the block first shifts item $j$
left by $k$ positions, so the reinsertion point becomes $j-k$ rather than $j$:

\begin{verbatim}
def move(D, i, k, j):
    items = []
    for t in range(k):
        items.append(D.delete_at(i))
    pos = j if j < i else j - k
    for t in range(k):
        D.insert_at(pos + t, items[t])
\end{verbatim}

This is again $2k$ calls to $O(\log n)$ operations, so $O(k\log n)$ total.
\end{problemparts}

\newpage
\problem  % Problem 3

Split the binder into three regions based on the two bookmarks: $L$ (pages before $A$),
$M$ (pages between $A$ and $B$), and $R$ (pages after $B$). Store each region in its own
dynamic array implemented as a circular buffer/deque (the same doubling/halving dynamic
array from lecture, but allowing $O(1)$ amortized push/pop at \emph{either} end, not just the
back). Each region also supports $O(1)$ worst-case random access by its local index. Maintain
two integers, $|L|$ and $|M|$, so that global index $i$ maps to local index $i$ in $L$ if
$i < |L|$, local index $i - |L|$ in $M$ if $|L| \le i < |L|+|M|$, and local index
$i - |L| - |M|$ in $R$ otherwise.

With ``in front of page $p$'' meaning ``immediately precedes $p$,'' the operations become:

\begin{itemize}
\item \textbf{build($X$):} Load every page from $X$ into a single deque (say $M$, with
$L$ and $R$ empty) using the standard $O(n)$ array-build. \emph{Worst-case} $O(n)$.

\item \textbf{place\_mark($i,m$):} Walk once through the current contents (in $O(n)$) and
redistribute pages into new $L$/$M$/$R$ deques consistent with the updated bookmark
positions, rebuilding each with the $O(n)$ array-build. \emph{Worst-case} $O(n)$.

\item \textbf{read\_page($i$):} Use $|L|,|M|$ to identify the region and local index, then do
one array access. \emph{Worst-case} $O(1)$.

\item \textbf{shift\_mark($m,d$):} Shifting $A$ moves one page between the $A$-end of $L$ and
the $A$-end of $M$ (\texttt{pop\_back}/\texttt{push\_front} or the reverse); shifting $B$
moves one page between the $B$-end of $M$ and the $B$-end of $R$ similarly. Update $|L|$ or
$|M|$ by $\pm 1$. \emph{Amortized} $O(1)$ (dynamic-array push/pop).

\item \textbf{move\_page($m$):} The page in front of $A$ is the last page of $L$; move it
(\texttt{pop\_back} on $L$, \texttt{push\_back} on $M$) to become the new last page of $M$,
i.e.\ in front of $B$. Symmetrically, the page in front of $B$ is the last page of $M$; move
it (\texttt{pop\_back} on $M$, \texttt{push\_back} on $L$) to become the new last page of $L$,
i.e.\ in front of $A$. Update $|L|,|M|$ accordingly. \emph{Amortized} $O(1)$.
\end{itemize}

Every operation only ever touches the two pages adjacent to a bookmark, which is exactly why
keeping $L$, $M$, $R$ as separate deques (rather than one big array) keeps
\texttt{shift\_mark} and \texttt{move\_page} from having to shift the rest of the binder.

\newpage
\problem  % Problem 4

\begin{problemparts}
\problempart % Problem 4a
Each of these only ever touches \texttt{head}, \texttt{tail}, and the one new/removed node,
so all four run in $O(1)$.

\textbf{insert\_first($x$):} Create a new node $y$ with $y.item = x$. If the list is empty
(\texttt{tail} is \texttt{None}), set $L.head = L.tail = y$. Otherwise, point the new node at
the old head ($y.next = L.head$), point the old head back at it ($L.head.prev = y$), then make
it the new head ($L.head = y$).

\textbf{insert\_last($x$):} Symmetric to \texttt{insert\_first}. Create a new node $y$. If the
list is empty, set $L.head = L.tail = y$. Otherwise, link the old tail forward to it
($L.tail.next = y$), link it back to the old tail ($y.prev = L.tail$), then make it the new
tail ($L.tail = y$).

\textbf{delete\_first():} If the list is empty, return \texttt{None}. Otherwise let
$x = L.head.item$. If $L.head.next$ exists, advance the head ($L.head = L.head.next$) and
clear its back-pointer ($L.head.prev = \texttt{None}$); if it doesn't (the list had exactly
one node), set $L.head = L.tail = \texttt{None}$. Return $x$.

\textbf{delete\_last():} Symmetric to \texttt{delete\_first}, using \texttt{tail}/\texttt{prev}
in place of \texttt{head}/\texttt{next}.

\problempart % Problem 4b
Since $x_1$ occurs before $x_2$ in $L$, everything from $x_1$ through $x_2$ becomes the new
list $L_2$, and $L$ keeps whatever came before $x_1$ and after $x_2$. All of this is just
pointer surgery around the two boundary nodes -- no traversal needed -- so it's $O(1)$:

\begin{enumerate}
\item Set $L_2.head = x_1$ and $L_2.tail = x_2$.
\item If $x_1$ has a predecessor in $L$ (i.e.\ $x_1 \ne L.head$), link that predecessor
directly to $x_2.next$, skipping the removed chain: $x_1.prev.next = x_2.next$. Otherwise
$x_1$ was the head of $L$, so $L.head = x_2.next$.
\item Symmetrically, if $x_2$ has a successor in $L$ (i.e.\ $x_2 \ne L.tail$), link it back to
$x_1.prev$: $x_2.next.prev = x_1.prev$. Otherwise $x_2$ was the tail of $L$, so
$L.tail = x_1.prev$.
\item Finally, cut $L_2$ loose from $L$ by clearing its outward-facing pointers:
$x_1.prev = \texttt{None}$ and $x_2.next = \texttt{None}$.
\end{enumerate}

\problempart % Problem 4c
Let $y = x.next$ (the node in $L_1$ that currently follows $x$; $y$ is \texttt{None} if
$x = L_1.tail$). If $L_2$ is empty, there's nothing to do. Otherwise, stitch $L_2$ in between
$x$ and $y$, then empty $L_2$ out -- again just pointer rewiring, so $O(1)$:

\begin{enumerate}
\item Link $x$ forward to $L_2$'s first node, and link that node back to $x$:
$x.next = L_2.head$, $\ L_2.head.prev = x$.
\item Link $L_2$'s last node forward to $y$: $L_2.tail.next = y$. If $y$ exists, link it back
to $L_2.tail$ ($y.prev = L_2.tail$); otherwise $x$ was $L_1$'s tail, so update
$L_1.tail = L_2.tail$.
\item Empty out $L_2$: set $L_2.head = L_2.tail = \texttt{None}$.
\end{enumerate}

\problempart Submit your implementation to {\small\url{alg.mit.edu}}.
\end{problemparts}

\end{problems}

\end{document}
